\documentclass[11pt]{article}
\usepackage[T1]{fontenc}
\usepackage{lmodern,amsmath,amssymb,amsthm}
\usepackage[margin=1in]{geometry}
\usepackage[hidelinks]{hyperref}
\setlength{\emergencystretch}{1em}
\newtheorem{theorem}{Theorem}
\newtheorem{lemma}[theorem]{Lemma}
\newtheorem{corollary}[theorem]{Corollary}
\newcommand{\Z}{\mathbb Z}
\newcommand{\Q}{\mathbb Q}
\newcommand{\F}{\mathbb F}
\newcommand{\aug}{\operatorname{aug}}
\newcommand{\rk}{\operatorname{rank}}
\newcommand{\len}{\operatorname{length}}
\newcommand{\Ph}[1]{\Phi_{p^{#1}}}
\title{Integral annihilators and the rank of a lattice with torsion cohomology}
\author{\small \href{mailto:admin@cosmosgeek.org}{admin@cosmosgeek.org}\qquad \href{mailto:b321014@iiit-bh.ac.in}{b321014@iiit-bh.ac.in}}
\date{Working paper for specialist review}
\begin{document}
\maketitle
\begin{abstract}
For a finite abelian $p$-group $G$, an element of order $p^4$ in $H^1(G,M)$
forces $\rk M-\rk M^G\ge p^4-1$ for every finite free $\Z_p$-lattice
$M$. We prove the bound by constructing an integral annihilator of
augmentation $p^3$ for every smaller invariant-free rational representation.
A cyclic augmentation lattice attains equality. We also give a feasibility
criterion for simultaneous elementary divisors of cyclic-group cohomology
with prescribed rational character multiplicities. In arithmetic the bound
applies to a Weil--Ch\^atelet class split by an abelian $p$-extension.
\end{abstract}

\section{The annihilator}
Fix a prime $p$. For an element $g$ in a group ring write
$N_q(g)=1+g+\cdots+g^{q-1}$ and
$\Ph{j}(g)=N_p(g^{p^{j-1}})$. For a subgroup $U$ let
$N_U=\sum_{u\in U}u$. A \emph{type of order $p^j$} means the Galois
orbit of a character of image order $p^j$. Over $\Q_p$ its irreducible
representation has dimension $p^{j-1}(p-1)$, since $\Ph{j}(X)$ is
irreducible. Killing one character from each occurring type kills the
entire rational representation.

\begin{theorem}\label{thm:main}
Let $m\in\{2,3,4\}$. If $G$ is a finite abelian $p$-group and
$V$ is a $\Q_p[G]$-representation with $V^G=0$ and
$\dim V<p^m-1$, there is $f\in\Z_p[G]$ satisfying
$fV=0$ and $\aug(f)=p^{m-1}$.
\end{theorem}

We denote this assertion by $A_m$. In the zero representation a scalar
suffices. One order-$p^m$ constituent leaves a complement of dimension
less than $p^{m-1}-1$. There can be at most one such constituent, counted
with multiplicity, and none of higher order. Choose $g$ on which the
character has order $p^m$. Multiplication by $\Ph{m}(g)$ treats this
constituent after applying $A_{m-1}$ to the complement. When $m=2$ the
complement has dimension below $p-1$ and hence is zero.

We record the reduction used in the other cases. Inflate the representation
to $(C_{p^{m-1}})^h$ for some $h\ge m-1$, and put $k=m-1$.
Annihilators push forward under the quotient map with augmentation preserved.
The dimension bound permits at most $p+\cdots+p^{m-1}$ distinct
order-$p$ types. In the dual of $\F_p^h$, fewer than
$1+p+\cdots+p^k$ specified projective points can be avoided by a
subspace of codimension $k$. Indeed, grow a disjoint subspace by one
dimension at a time. Before the final extension the quotient has at
least $1+p+\cdots+p^k$ projective directions, and each specified point
excludes at most one. Taking the annihilator of the resulting subspace
gives a $k$-plane on which every specified character is nontrivial.

There are at most $1+p+\cdots+p^{m-2}<p^{m-1}$ higher types.
For an order-$p^j$ character whose primitive defining linear form
vanishes modulo $p$ on this plane, a random lift of each basis vector
lands in its kernel with probability $p^{-(j-1)}$. The $k$ lifts are
independent. A union bound therefore gives a lift on which no type is
trivial as a character of the whole subgroup. The lifted basis
generates $(C_{p^{m-1}})^k$. Its annihilator works in the ambient group
ring. We may hence work on this homocyclic subgroup whenever no
order-$p^m$ constituent occurs.

\section{The first two cases}
For $A_2$ with no order-$p^2$ constituent there are at most $p$
order-$p$ types. At most $p$ proper linear hyperplanes cannot cover
a nonzero finite vector space over $\F_p$, as they all contain zero.
Choose $g$ outside their kernels and take $\Ph{1}(g)$.

For $A_3$ the preceding reductions leave $G=(C_{p^2})^2$.
Let $a,b$ count distinct types of orders $p,p^2$. Then
$a\le p+1$ and $a+pb\le p^2+p$. If $a=0$, use $N_{pG}$,
with augmentation $p^2$; if $b=0$, take $\Ph{1}(x)\Ph{1}(y)$
for a basis. The case $b=p+1$ requires $a=0$.

If $1\le b<p$, choose a basis vector $x$ primitive on all high types
(i.e. its character values have order $p^2$), and extend it by $y$.
Normalize representatives so their $x$-values equal
$\zeta=\zeta_{p^2}$; their $y$-values $\alpha_i$ are distinct.
Set $O=\Z_p[\zeta]$ and $\pi=\zeta-1$. On dividing $T^n$ by
$\prod_i(T-(\alpha_i-1))$, the coefficient of $T^i$ in the remainder
lies in $\pi^{n-i}O$ for $i<b$ and $n\ge b$. This follows inductively
from the elementary symmetric coefficients. The coefficients of
$\Ph{1}(1+T)$ below degree $p-1$ are divisible by $p$, and the
remainder of $T^{p-1}$ is divisible by $\pi$. Thus an integral
polynomial in $Y$ takes values $-\Ph{1}(\alpha_i)/\pi$ at $Y=\alpha_i$.
Lift its coefficients from $O$ through $X\mapsto\zeta$, obtaining $A$.
The element
\[
 \Ph{1}(x)\Ph{1}(y)+(x^p-1)A(x,y)
\]
annihilates the indicated high types since $\Ph{1}(\zeta)=(\zeta^p-1)/\pi$.
It also annihilates every low type and has augmentation $p^2$.

When $b=p$, $a\le p$. Take a nonzero reduced vector $v$ outside the
low kernels. A high type reducing trivially on $v$ forbids one affine
line among the lifts $g=g_0+ph$, $h\in\F_p^2$. If these lines fail to
cover, $N_{p^2}(g)$ suffices. A cover by exactly $p$ affine lines
requires all of them to be parallel and distinct. Thus every high
type has the same reduced kernel $L$, and $v\in L$. If any allowed
$v$ is outside $L$, use that vector. Otherwise the low types are exactly
the $p$ types with kernels different from $L$. Choose basis vectors with
$y$ reducing into $L$ nontrivially and $x$ outside it. Then
$\Ph{2}(x)\Ph{1}(y)$ works. This proves $A_3$.

\section{An interpolation in rank three}
\begin{lemma}\label{lem:interp}
In $(C_{p^2})^3$ with basis $x,y,z$, let fewer than $p^2$ selected high
types be primitive on $x$. There is an annihilator of augmentation $p^3$
for these types, all nontrivial order-$p$ types and any further types
taking a nontrivial order-$p$ value on $x$.
\end{lemma}
\begin{proof}
Normalize the selected $x$-values to $\zeta_{p^2}$. The joint
$O=\Z_p[\zeta_{p^2}]$-evaluation algebra $B$ generated by $y,z$ in
$O^{b_0}$ has rank $b_0$, the number of selected types. Over the
fraction field, interpolation at distinct pairs of values separates the
coordinates. Put $\pi=\zeta_{p^2}-1$. The quotient
$D=B/(\pi,y^p-1,z^p-1)$ has dimension at most $b_0<p^2$ over $\F_p$.
It is a quotient of $\F_p[S,T]/(S^p,T^p)$ via $S=y-1$, $T=z-1$.
Any nonzero ideal in this local ring contains its one-dimensional socle
generated by $S^{p-1}T^{p-1}$: take the last nonzero power of the maximal
ideal acting on the given ideal. Therefore $\Ph{1}(y)\Ph{1}(z)$ vanishes
in $D$. Lifting the membership identity to the group ring yields
\[
 \Ph{1}(x)\big[\Ph{1}(y)\Ph{1}(z)
 -(x-1)A_0-(y^p-1)A_1-(z^p-1)A_2\big].
\]
It has the required zeros and augmentation. If $b_0=0$ take the product
of the three cyclotomic factors.
\end{proof}

Consider $A_4$ first for groups of exponent at most $p^2$. We may
restrict to $(C_{p^2})^3$. Indeed the preceding $k=3$ avoidance
argument applies after inflation at exponent $p^2$; there are at most
$p^2+p+1<p^3$ high types. On this rank-three group let $a,b$ count
low and high types. Put $P=p^2+p+1$. The budgets are
$a\le P$ and $a+pb\le p^3+p^2+p$.

For a uniformly chosen nonzero vector mod $p$, the expected number $c$
of high reduced kernels containing it is $b(p+1)/P$. For $b<p^2$
choose $c<p$ and lift the vector $x$ to avoid the at most $p-1$ affine
hyperplanes which would give trivial values. The bad types then have
nontrivial order-$p$ values on $x$. Apply Lemma~\ref{lem:interp} to
the other $b-c$ types.

For $b=p^2+j\ge p^2$, the case $a=0$ is handled by $N_{pG}$.
Otherwise $j\le p$, and averaging gives $c\le p$.
If $c\le j$, the multiplicity-free low and bad support has dimension
$(p-1)(a+pc)\le(p-1)(p^2+p)<p^3-1$. Use $A_3$ on that support and
multiply by $\Ph{2}(x)$ to kill the good types. If $c>j$ then
$b-c<p^2$, so Lemma~\ref{lem:interp} applies after avoiding trivial
values on the $c$ bad types. The only remaining possibility is
$c=p>j$ with a cover by $p$ parallel affine hyperplanes. These $p$
bad types share a reduced functional up to a scalar, say $\alpha$.
Condition instead on $\alpha(x)\ne0$. Each other high kernel contains
at most proportion $1/p$ of these choices, giving mean
$(b-p)/p=p-1+j/p<p$. We return to one of the earlier cases.

\section{Two integral lemmas}
\begin{lemma}\label{lem:norm}
Let $P_0=(C_p)^h$, $h\ge1$, $n=p^h$, and let $D$ be a finite module
over $R=(\Z/p^2)[P_0]$. If
$\len_{\Z_p}D<n$ and $\dim_{\F_p}pD<p^{h-1}$, then $N_{P_0}D=0$.
\end{lemma}
\begin{proof}
Put $A=\F_p[P_0]=\F_p[T_1,\ldots,T_h]/(T_1^p,\ldots,T_h^p)$, where
$T_i=g_i-1$ for a basis of $P_0$. Suppose multiplication by $f\in A$
has rank less than $p^{h-1}$. If the first nonzero homogeneous part
had degree $d\le p-1$, choose a leading monomial $T^e$ in that part.
Its multiples by $T^b$ with $0\le b_i<p-e_i$ give independent multiples
of $f$, by their distinct leading terms in the associated graded ring.
Their number is $\prod_i(p-e_i)\ge p^{h-1}(p-d)$ (concentrating
the exponents in one variable minimizes the product). Thus $f$ lies
in the $p$th power of the augmentation ideal.

Now take $F\in\Z_p[P_0]$ with augmentation of valuation one and
reduction of multiplication rank less than $p^{h-1}$. In its $T$-basis
expansion, coefficients of degrees below $p$ are divisible by $p$.
At any character of $P_0$ every nonconstant term has
$(\zeta_p-1)$-valuation strictly greater than $p-1$, whereas the constant
term has valuation exactly $p-1$. Consequently every character value
of $F$ has $p$-valuation one. Multiplication by $F$ has determinant
of valuation $n$ by the character decomposition. If its Smith
exponents are $a_i$, its image in $R$ has length
$\sum_i\max(2-a_i,0)\ge 2n-\sum_i a_i=n$.

Suppose $N_{P_0}d\ne0$ and replace $D$ by $Rd$. In characteristic $p$,
a vector on which the norm is nonzero generates a free rank-one copy
of $A$, since every nonzero ideal of $A$ contains the norm (the
socle). Hence $N_{P_0}$ kills $pD$ and $N_{P_0}d$ has order $p$.
Choose a functional $\lambda:D\to\Z/p^2$ with
$\lambda(N_{P_0}d)=p$; this follows by decomposing the underlying
abelian group into cyclic factors. The contragredient dual $D^*$
embeds into $R$ by $\lambda\mapsto f=\sum_g\lambda(gd)g$,
since the $gd$ generate $D$. This map is equivariant, and
$\aug(f)=p$. The rank of multiplication by $f$ modulo $p$ is
$\dim_{\F_p}pRf\le\dim_{\F_p}pD^*=\dim_{\F_p}pD<p^{h-1}$.
For the equality, multiplication by $p$ identifies $R/pR$ with $pR$.
The determinant bound applied to a lift of $f$ gives $\len(Rf)\ge n$,
contrary to $\len D^*=\len D<n$.
\end{proof}

\begin{lemma}\label{lem:high}
In $G=(C_{p^3})^3$ let $c$ distinct order-$p^3$ types be specified,
where $1\le c<p$. For any $u,v\in G$, either of the elements
\[
 f_0=N_{p^2}(v)N_p(u),\qquad f_0=N_{p^2}(v)N_p(u^p)
\]
can be altered by an augmentation-zero correction that kills all $c$
types and leaves unchanged its values on every type of order at most $p^2$.
\end{lemma}
\begin{proof}
Choose $x$ primitive on every specified type and normalize $x$ to
$\zeta=\zeta_{p^3}$. Let $B\subset O^c$ be the joint evaluation
algebra, with $O=\Z_p[\zeta]$ and $\pi=\zeta-1$. In $B$ put
$I=(u-1,v-1)B$. Its coordinates lie in $\pi O$. For $a\in I$,
Cayley--Hamilton for the diagonal coordinates gives
$a^c\in\sum_{j=1}^c\pi^j I^{c-j}$. Polarization applies as $c!$ is
invertible and yields the same containment for $I^c$. Induction gives
\[
 I^d\subseteq\sum_{i=0}^{c-1}\pi^{d-i}I^i
 \subseteq\pi^{d-c+1} B\qquad(d\ge c-1).
\]
Modulo $p$ each product in the assertion consists of its top monomial
in $u-1,v-1$, of degree respectively $p^2+p-2$ and $2p^2-p-1$.
Both degrees satisfy $d-c+1\ge p^2$. As $v_\pi(p)=p^2(p-1)$,
$f_0\in\pi^{p^2}B=(\zeta^{p^2}-1)B$. Lift the quotient to
$A\in\Z_p[G]$ and use $f_0-(x^{p^2}-1)A$.
\end{proof}

\section{The highest remaining character types}
It remains to prove $A_4$ for $(C_{p^3})^3$ with at least one
order-$p^3$ type. Let $C$ count those constituents with multiplicity,
$c$ their distinct types, and let $a,b$ count distinct types of orders
$p,p^2$. Then
\[
 a+pb+p^2C\le p^3+p^2+p,\qquad C\le p+1.
\]
If $C=p$, the entire lower complement has dimension below $p^3-1$.
Choose $x$ primitive on the high types; multiply an $A_3$ annihilator
by $\Ph{3}(x)$. If $C=p+1$, the lower complement has dimension below
$p^2-1$. Two vectors, drawn independently modulo $p$, simultaneously
miss a fixed reduced kernel in the sense that at least one is outside,
except with probability $1/p^2$. Since $(p+1)/p^2<1$, two suffice for
all high types. Use the product of the two $\Ph{3}$ factors and an
$A_2$ annihilator. We can now assume $1\le c\le C<p$.

If $a<P=p^2+p+1$, some missing low type $\alpha$ has fewer than $p^2$
middle types whose reduced functionals are proportional to $\alpha$.
Otherwise, putting $t=P-a\ge1$, we would have
$a+pb+p^2C\ge P+(p^3-1)t+p^2C>p^3+p^2+p$.
All present low types are nontrivial on $\ker\alpha$.
In $G/p^2G$ lift a basis of this plane: each of the fewer than $p^2$
obstructing middle types forbids proportion $p^{-2}$ of the pairs of
lifts. The resulting subgroup $H=(C_{p^2})^2$ has no fixed vector on
the lower part, whose dimension is below $|H|-1$. Consequently some
nontrivial rational $H$-type $\psi$ is absent: the direct sum of one
copy of each nontrivial rational type has dimension $|H|-1$.

Write $K_0=\ker\psi$ and let $g$ generate $H/K_0$. If $\psi$ has
order $p^j$, then $N_{K_0}N_{p^{j-1}}(g)$ has augmentation $p^3$.
Its Fourier support consists of the trivial type and the type of
$\psi$: the norm first forces triviality on $K_0$, and the geometric
sum kills the remaining nontrivial lower orders. An adapted basis
$u,v$ yields $N_{p^2}(v)N_p(u)$ when $j=2$ and
$N_{p^2}(v)N_p(u^p)$ when $j=1$. Lift $u,v$ to $G$ and apply
Lemma~\ref{lem:high} to kill all high types.

Finally suppose $a=P$. Then $b\le p(p-C)-1\le p(p-c)-1$.
The complement $\Omega$ of the high kernel planes in $G/pG$ has size
at least $(p-c)p^2$. Each middle kernel contains at most $p^2$ points
of $\Omega$, so some $x_0\in\Omega$ lies in at most $p-1$ middle
kernels. Lift it to $x$ to avoid their trivial-value hyperplanes. The
bad middle types then take nontrivial order-$p$ values on $x$ and the
other $b_0$ middle types take primitive order-$p^2$ values. All high
types are primitive on $x$. Complete $x$ to a basis $x,y,z$.

Evaluate jointly on the $b_0$ good middle types and the $c$ high
types, normalizing $x$ to the respective primitive roots. The image
$B$ of $\Z_p[G]$ is a full lattice in
$O_2^{b_0}\times O_3^c$ for $O_j=\Z_p[\zeta_{p^j}]$, by rational
semisimplicity. Multiplication by $x-1$ has determinant valuation
$b_0+c$. Thus $\bar B=B/(x-1)B$ has length $b_0+c<p^2$.
Let $B_2$ be the middle projection and $J$ its kernel. Since $x-1$
is injective on $B_2$, we have an exact sequence
\[
 0\longrightarrow J/(x-1)J\longrightarrow\bar B
 \longrightarrow B_2/(x-1)B_2\longrightarrow0.
\]
Both ends are killed by $p$, since their respective cyclotomic
polynomials evaluate to $p$ at $1$. The first term has length $c$ by
the determinant on the high coordinates. Hence $p^2\bar B=0$ and
$\dim_{\F_p}p\bar B\le c<p$.

The quotient $D=B/(x-1,y^p-1,z^p-1)$ is a module over
$(\Z/p^2)[(C_p)^2]$, of length less than $p^2$ with
$\dim_{\F_p}pD<p$. Lemma~\ref{lem:norm} gives
$\Ph{1}(y)\Ph{1}(z)=0$ in $D$. Lift this identity exactly as in
Lemma~\ref{lem:interp}, with $x-1$ replacing $\pi$. The bracket there
kills the good middle and high types, and $\Ph{1}(x)$ kills the bad
middle types. On a nontrivial low type with $x=1$, at least one of
$\Ph{1}(y),\Ph{1}(z)$ vanishes and all correction terms vanish.
The element has augmentation $p^3$, completing Theorem~\ref{thm:main}.

\section{Cohomological and arithmetic consequences}
\begin{corollary}
If $M$ is a finite free $\Z_p$-module with an action of a finite
abelian $p$-group $G$ and $H^1(G,M)$ contains an element of order
$p^4$, then $\rk M-\rk M^G\ge p^4-1$. The bound is attained.
\end{corollary}
\begin{proof}
Let $M_0=\ker(N_G:M\to M)$. Its rationalization has no invariants,
and its rank is the asserted rank difference. Since $N_GM$ is a
trivial torsion-free module, the map $H^1(G,M_0)\to H^1(G,M)$ is
surjective. A cocycle $z$ on an abelian group satisfies
$(h-1)z(g)=(g-1)z(h)$; hence $h$ acts trivially on its cohomology
class. Every group-ring element acts on $H^1$ by its augmentation.
If the moving rank were smaller than $p^4-1$, Theorem~\ref{thm:main}
would give an integral element killing $M_0$ and acting by $p^3$ on
$H^1$. This contradicts the assumed order. For $G=C_{p^4}$ take the
augmentation ideal in $\Z_p[G]$. It has no invariants and rank
$p^4-1$. The regular-module augmentation sequence gives
$H^1=\Z_p/p^4\Z_p$.
\end{proof}

If $F/k$ is a finite abelian $p$-extension of number fields, $E/k$ is
elliptic and $E(F)[p]=0$, the same bound applies to any order-$p^4$
Weil--Ch\^atelet class split by $F$. Indeed inflation--restriction
identifies the kernel of restriction with $H^1(G,E(F))$.
Prime-to-$p$ point torsion has vanishing higher cohomology by averaging.
Completion of the free lattice at $p$ preserves this $p$-primary
cohomology (use the finite-group cochain complex and flatness).
The resulting moving rank is $\rk E(F)-\rk E(k)$. More generally,
if the $p$-primary point torsion has exponent dividing $p^b$, the kernel
of passage to the free lattice has exponent dividing $p^b$, so a class
of order $p^{b+4}$ split by $F$ gives the same bound.

\section{Cyclic lattices with prescribed multiplicities}
Let $G=\langle s\rangle=C_{p^n}$, $L=\ker N_G$ in $M$, and write
\[
 L\otimes\Q_p=\bigoplus_{j=1}^n\Q_p(\zeta_{p^j})^{c_j}.
\]
For $A=\bigoplus_i\Z/p^{a_i}$ with $a_i\ge1$, a necessary condition for
$A\hookrightarrow H^1(G,M)$ is, for every $J\subseteq\{1,\ldots,n\}$,
\begin{equation}\label{eq:flow}
 \sum_i(a_i-|J|)_+\ \le\sum_{j\notin J}c_j .
\end{equation}
These inequalities are sufficient for some invariant-free lattice of
the specified rational type.

Indeed $H^1(G,M)=L/(s-1)M$ is a quotient of $C=L/(s-1)L$. Intersect
$L$ with the rational constituents indexed by $J$ to obtain a saturated
$L_J$. Taking $(s-1)$-coinvariants preserves the resulting short exact
sequence since the quotient has no fixed vector. The image of
$(L_J)_G$ in $C$ is killed by $p^{|J|}$: the polynomial
$\prod_{j\in J}\Ph{j}(X)$ annihilates $L_J$ and takes this value at $1$.
Multiplication by $p^{|J|}$ therefore induces a surjection
$(L/L_J)_G\twoheadrightarrow p^{|J|}C$. The source has order
$p^{\sum_{j\notin J}c_j}$ by the determinant of $s-1$. This bounds the
order of $p^{|J|}A$ and proves \eqref{eq:flow}.

Conversely assign row $i$ to $a_i$ distinct levels, with column
capacities $c_j$. The integral max-flow theorem gives such an
assignment exactly under \eqref{eq:flow}: take source-to-row
capacities $a_i$, row-to-level capacities one, and level-to-sink
capacities $c_j$. For levels $J$ on the sink side of a cut,
minimizing over the row placements gives cut capacity
$\sum_i\min(a_i,|J|)+\sum_{j\notin J}c_j$.
For the levels $S_i$ assigned to row $i$, use
$\Z_p[X]/(\prod_{j\in S_i}\Ph{j}(X))$ with $s=X$. Its coinvariants
are cyclic of order $p^{a_i}$. Fill unused column capacity by
individual cyclotomic lattices. This realizes $A$ as a subgroup of $H^1$.

In particular $\rk L\ge\sum_i(p^{a_i}-1)$: each row costs at least
$\sum_{j=1}^{a_i}\varphi(p^j)=p^{a_i}-1$. Using the first $a_i$
levels for each row attains equality at the level of integral lattices.
This numerical cyclic bound is due to Edixhoven
\cite[Lemma 4.5, Remark 4.6 and Corollary 4.7]{Ed}; see also
Bosch--Xarles \cite[Proposition 5.11]{BX}. The cut criterion keeps the
individual multiplicities as input. For a cyclic $p$-extension and
$E(F)[p]=0$ the same inflation--restriction argument bounds the rank
gain when several independent classes split simultaneously.

\paragraph{A restriction-of-scalars check on integral splitting.}
Rational rank jumps do not in general determine the integral lattice.
For an odd $p$, take a $\Z_p$-extension $F_\infty/F$ with first
degree-$p$ layer $L$, and an elliptic curve $D/F$ of positive rank.
Put $B=\operatorname{Res}_{L/F}(D_L)$ and $r=\rk D(L)>0$. Then
$B(F)=D(L)$ and $B(L)\simeq D(L)^p$. In the coordinates coming from
$L\otimes_F L=\prod_{g\in G}L$, the action is
$(s x)_g=s(x_{s^{-1}g})$. Replacing $x_g$ by $g^{-1}x_g$ leaves only
the permutation of factors. Thus
\[
 T_p(B(L)\otimes\Q_p/\Z_p)\simeq\Z_p[C_p]^r .
\]
The integral direct sum asserted in the display of Proposition 3.8 of
the cited version of \cite{Lim} instead gives
$(\Z_p\oplus\Z_p[\zeta_p])^r$, as the multiplicities there are
$e_0=e_1=r$. The regular module has zero $H^1$, whereas the displayed
direct sum has $H^1=(\F_p)^r$. Evaluation at $1$ and $\zeta_p$
identifies the regular lattice with the pairs whose residues agree
modulo $\zeta_p-1$; the index is $p$ per summand. Both rationalized
representations are the same.

\section{Context}
Edixhoven's cyclic bound bears on a question of Xarles on noncyclic
lattice cohomology. Theorem~\ref{thm:main} treats one order-$p^4$
class for abelian groups of arbitrary $p$-power order. Adem's results
on maximal cohomological exponent \cite{Adem} concern the case of
exponent $|G|$, which can exceed $p^4$. Establishing literature priority
for the exact noncyclic statement requires further specialist comparison.
The arithmetic corollary concerns growth in the splitting field and
makes no comparison with the complex analytic rank.

\begin{thebibliography}{9}
\bibitem{Ed} B. Edixhoven, \emph{On the prime-to-$p$ part of the
groups of connected components of N\'eron models}, Compositio Math.
\textbf{97} (1995), 29--49.
\href{https://www.numdam.org/item/CM_1995__97_1-2_29_0/}{Article}.
\bibitem{BX} S. Bosch and X. Xarles,
\emph{Component groups of N\'eron models via rigid uniformization},
Math. Ann. \textbf{306} (1996), 459--486.
\href{https://mat.uab.es/~xarles/components.pdf}{Article}.
\bibitem{Adem} A. Adem, \emph{Cohomological exponents of
$\mathbb ZG$-lattices}, J. Pure Appl. Algebra (1989).
\href{https://doi.org/10.1016/0022-4049(89)90048-0}{Article}.
\bibitem{Lim} M. F. Lim, \emph{On fine Mordell--Weil groups over
$\mathbb Z_p$-extensions of an imaginary quadratic field},
\href{https://arxiv.org/pdf/2308.04096v3}{arXiv:2308.04096v3},
Proposition~3.8.
\end{thebibliography}
\end{document}
